We measure how label noise and ridge regularisation change the height and position of the test-error peak of random-ReLU-features regression, and whether a tuned ridge removes it. The study is small and CPU-only: $n=200$ training points, 24 widths from $N/n=0.05$ to $25$, a synthetic single-index task and sklearn digits regressed on one-hot labels, four label-noise levels each, and 5 seeds. For minimum-norm interpolation the peak sits at $N/n=1$ in every one of the 40 runs and rises with label noise (one-sided paired Wilcoxon $p=0.031$ on both datasets, the smallest value possible with 5 seeds); it is, however, already very large without noise (mean peak test MSE above the best error by a factor of $10^3$ to $10^5$). A fixed ridge lowers the peak by orders of magnitude, but our registered hypothesis that the peak decreases monotonically in the ridge is refuted: the peak is smallest at $\lambda=10^{-3}$ or $10^{-2}$ and larger again at $\lambda\ge10^{-1}$, where the curve is flat. A ridge tuned per width on a validation set leaves a residual hump below 5\% of the best error on average on all 8 tasks (largest task mean 0.033), but individual seeds on three tasks exceed 5\%, and the registered paired test against minimum-norm interpolation on the hump statistic is not significant at 5 seeds because of its heavy tail. A single global ridge removes the hump as well or better than per-width tuning in our data; leave-one-out and small-validation tuning do not on the noisiest synthetic task. All claims are for this finite-size, small-scale setting.
